\chapter{Build Something That Can Outlive You}
\chaptermark{Build Something That Can Outlive You}

An older businessman approached on a London street makes the sale of a company
sound like the decisive act. Starting the business creates work and cash-flow
pressure, he says; selling it creates wealth. When pressed for a method, he
advises a founder to identify likely buyers at the beginning, learn the trade,
build a client base, grow turnover, and make the company useful to one of those
buyers.

The advice is memorable because it turns an exit from a distant event into a
design question. It is incomplete because sale is not the only honorable end of
a company, and a buyer's preference is not the same as a customer's need. Some
businesses should remain independent. Some should pass to family, employees, a
management team, a cooperative, a trust, or another qualified steward. Some
should close cleanly rather than transfer work, data, risk, or promises to an
owner unable to honor them.

The deeper test is useful in every case: can the value continue when control
changes? A transferable company can keep its promises without hiding its
obligations or depending on one person's presence. It knows what it owns, what
it owes, why customers stay, how decisions are made, who can lead next, and
which commitments must survive any transaction.

Saleability is therefore not the purpose of the company. It is a demanding test
of whether the company has become more than a job held together by its founder.

\section{Design for the next steward, not merely the next buyer}

The London businessman's four-year buyer horizon is his heuristic, not a law.
Its best use is to force specificity. If a larger firm might one day want the
company, what would it actually be unable or unwilling to build for itself? The
answer might be trusted customer access, a capable team, a license, a body of
knowledge, a repeatable service, a difficult location, a product with retained
users, a reputation in a narrow market, or a combination that works only when
kept together.

Now ask the same question without assuming a sale. What would a successor need
to preserve the customer outcome? What would employees need in order to work
without guessing at the founder's intent? What would suppliers and lenders need
to continue? Which permits, consents, contracts, and relationships survive a
change of control? Which promises depend on personal trust and must be renewed
rather than assigned?

These questions distinguish four things that are often collapsed into one:

\begin{description}[leftmargin=0em,style=unboxed,itemsep=0.1em,
  topsep=0.2em,parsep=0em,labelsep=0.55em]
\item[Continuity] the work can continue through an absence, illness, or ordinary
  leadership change.
\item[Succession] a named person or group can assume authority, judgment, and
  accountability over time.
\item[Transfer] rights and obligations can lawfully move to another owner with
  limited loss of customers, people, permission, and capability.
\item[Sale] a negotiated transfer produces consideration for specified rights,
  assets, or shares under agreed conditions.
\end{description}

A company may achieve continuity without being saleable, or be legally
saleable without being able to preserve its value. A professional practice can
have excellent records and a trained successor while remaining subject to
licensing and patient choice. A software company can transfer its shares while
its largest customer is free to cancel. A family company can produce reliable
cash and still have no willing or capable next steward.

Design begins with the destination. Write the acceptable paths, the prohibited
ones, and the conditions under which each becomes responsible. A buyer list is
useful only after this stewardship list exists.

\section{Remove dependence without erasing judgment}

The founder of a fitness application, reflecting on a reported sale, identifies
her own continued attachment to the company as the mistake. She wished she had
removed herself operationally before closing; instead of celebrating and
relaxing, she remained a critical dependency. The point is not that founders
should disappear. It is that a transaction cannot transfer a capability that
still exists only inside one person.

Chapter 11 built an absence-tolerant operating system. Transfer asks for a
harder test. A temporary backup can follow a procedure; a successor must also
understand why the procedure exists, when it no longer applies, and who bears
the consequence of changing it. Documentation must therefore preserve more
than steps. It should expose:

\begin{enumerate}[itemsep=0.12em]
\item the customer outcome and standard;
\item the evidence behind important decisions;
\item the limits within which a role may act;
\item the exceptions that require judgment or professional review;
\item the relationships and permissions that cannot be compelled;
\item the unresolved risks, disputes, and experiments; and
\item the history of failures, repairs, and promises still in force.
\end{enumerate}

One solar-company founder describes the transition from a business that relies
on the chief executive to one that relies on documented functions, people, and
systems. A private-equity operator similarly attributes his ability to oversee
many companies to common data, empowered people, and a repeatable process.
Those are useful transfer mechanisms, but neither uniformity nor a large manual
proves that the system is safe. Local regulation, customer context, tacit craft,
and dissent can be destroyed by forcing every company through one template.

The transfer test is practical. Let the prospective successor chair a real
operating review. Let the founder remain silent unless a stated safety or legal
boundary is reached. Observe which decisions stall, which facts are absent,
which relationships call only the founder, and which team members hold
responsibility without authority. Repair those failures while there is still
time to teach, not during a closing or crisis.

\section{Know what the headline price means}

A deal announcement can be impressive while revealing little about what any
owner receives. In one interview, a hospitality operator corrects the language
around a multibillion-dollar transaction from equity value to enterprise value.
In another, an energy-company seller reports a range because the transaction
included contingencies. These distinctions are not footnotes. They are the
difference between the value assigned to an operating business and cash that is
available to a particular person.

A simplified seller bridge is:

\begin{align}
P_{\mathrm{pre\mbox{-}tax}} \approx{}& EV + C - D - L - F \notag\\
&{}\pm W - H,
\end{align}

where \(EV\) is enterprise value, \(C\) is cash included for the seller,
\(D\) is debt repaid or assumed, \(L\) is debt-like and other negotiated
liability, \(F\) is transaction cost, \(W\) is the working-capital adjustment,
and \(H\) is money held back, placed in escrow, or dependent on later events.
The amount attributable to one owner then depends on the securities, ownership
percentage, preferences, options, taxes, indemnities, and earn-out terms that
apply to that owner. Jurisdiction and transaction structure change the result;
qualified legal, tax, and financial review is necessary.

Even a carefully stated enterprise value is not a promise that the price will
be paid. A buyer may condition payment on retention, performance, regulatory
approval, financing, representations, or a later milestone. A founder who must
remain for years to earn the headline price has sold a different object from a
founder who receives cash at closing and leaves.

The same discipline applies before valuation. Recurring revenue can make a
company easier to understand because retention provides evidence about future
demand. The London health-technology founder explains why institutions may
value a subscription stream differently from repeated one-off sales. A wealth
manager in another interview also favors retained, non-transactional client
economics. Neither account means that every subscription is durable. A buyer
will ask whether renewal is voluntary, whether customers receive continuing
value, what service and recovery cost, how cohorts behave, and whether the
contract survives transfer. Lock-in can make revenue look stable just before
reputation or regulation breaks it.

\section{Let trust travel without selling a person}

Robert Miller reports closing two agencies while continuing to build a public
reputation around digital marketing and client acquisition. He attributes the
faster launch of later businesses to that accumulated reputation. The offers
and market mechanisms changed; recognition of his claimed capability traveled.

Portable trust is real, but a personal brand is not immortal and follower count
is not trust. It depends on platforms, records, conduct, continued relevance,
and the truth of the claims attached to the name. It can also intensify
key-person risk. If every product borrows credibility from one person, one
misleading claim, private failure, health event, or departure can damage every
offer at once.

A Beverly Hills jeweler describes honor, values, and keeping one's word as the
assets from which he could rebuild if the visible business disappeared. That
claim identifies the right substance: reputation is the remembered pattern
between promise and conduct. It becomes transferable only when the company,
not merely its public face, repeatedly keeps the promise.

Build three connected but distinct forms of trust:

\begin{description}[leftmargin=2.8em,style=nextline,before=\raggedright]
\item[Personal trust] follows a person's demonstrated judgment and conduct.
\item[Product trust] follows evidence that a particular offer performs safely
  and reliably for the people it serves.
\item[Institutional trust] follows the organization's standards and records.
  It also follows its remedies, governance, and willingness to correct itself
  across personnel.
\end{description}

Do not try to transfer trust through a new logo and an announcement. Preserve
the evidence behind it: product history, quality measures, complaint and repair
records, warranties, customer permissions, truthful claims, professional
standards, community commitments, and the people who know how to act when the
promise is under pressure. Inform affected people honestly about a change of
control and give them the choices law and fairness require.

\section{Decide what should not be sold}

An exit can reward years of useful work, release concentrated household risk,
give a company resources it could not obtain alone, and place it with a better
steward. It can also move power to an owner whose economics conflict with the
reason customers, employees, or communities trusted the company.

Before optimizing for price, identify the non-price terms. They may include
continued service to vulnerable customers, worker safety, earned compensation,
privacy, product quality, location commitments, environmental repair,
professional independence, access to essential records, treatment of suppliers,
or preservation of a mission. Some can be protected by structure and contract;
some depend on choosing a credible steward; some cannot be protected well enough
to justify the sale.

Then distinguish a preference from an obligation. A founder's desire that the
brand remain unchanged is not automatically superior to the team's ability to
adapt. A worker's earned pay, a customer's privacy right, a lender's claim, a
regulatory duty, or a remediation obligation cannot be dismissed as sentimental
friction. Transfer does not cleanse the history of the asset.

The strongest offer is not necessarily the highest number. Compare certainty,
timing, financing, contingencies, governance during transition, employment and
customer effects, successor capability, legal exposure, and what remains at
risk after closing. A responsible refusal is sometimes more valuable than a
fragile headline.

\section{Run the transferability review}

Choose the company, practice, product, or income stream most important to the
life you are building. Review it as if control had to change within twelve
months, even if no sale is planned.

\begin{enumerate}[itemsep=0.14em]
\item \textbf{Customer promise.} State the outcome, evidence, quality standard,
  open complaint, and consequence of failure.
\item \textbf{Demand quality.} Measure customer and channel concentration,
  voluntary retention, cohort contribution, cancellation rights, and the
  founder's personal share of sales.
\item \textbf{Economics.} Reconcile revenue, margin, cash collection, working
  capital, capital expenditure, recurring service cost, and normalized owner
  compensation.
\item \textbf{Rights and contracts.} Verify ownership, license, consent,
  assignment, change-of-control, exclusivity, security, and termination terms.
\item \textbf{Records and controls.} Test whether accounts, taxes, payroll,
  decisions, incidents, approvals, access, and backups are complete and usable.
\item \textbf{People and succession.} Name the successor, essential knowledge,
  decision authority, retention risk, development plan, and lawful treatment
  of every affected worker.
\item \textbf{Operating continuity.} Rehearse delivery, collection, recovery,
  escalation, and learning without the founder's intervention.
\item \textbf{Reputation.} Separate personal, product, and institutional trust;
  list every public claim or unresolved harm that could travel with the name.
\item \textbf{Obligations.} Surface debt, guarantees, liens, litigation, tax,
  privacy, environmental, warranty, safety, and community commitments.
\item \textbf{Stewardship choice.} Compare continued ownership, family or
  management succession, employee ownership, outside sale, partial transfer,
  partnership, and orderly closure against the life and duties at stake.
\end{enumerate}

The review should produce named repairs, owners, evidence, and dates. Repeat it
after a material contract, product, financing, leadership, or regulatory change.
If the company becomes easier to transfer, it should also become safer to own:
less dependent on one body, less opaque about cash, more honest about risk, and
more capable of keeping its promises through change.

That is the end of Part III's argument. Skill became an owned claim; the claim
became a system; the system acquired people, cash needs, rights, and a future
beyond its founder. Part IV asks what capital can do to that machine. Used with
discipline, it can make the company resilient. Mispriced or misunderstood, it
can turn every advantage into another way to fail.
